% LSPV-1 REPORT::LDELEAF1::2026-10-01T20:38:00-07:00
% !TEX program = xelatex
\documentclass[11pt]{article}
\usepackage[margin=0.88in]{geometry}
\usepackage{amsmath,amssymb}
\begin{document}
\section*{LSPV-1: source-typed vMS LDE leaf and Hadamard exterior test}
\textbf{STATUS: OPEN.} P versus NP remains open. No unrestricted SAT lower bound, nonrelativizing hinge, or nonalgebrizing hinge is proved.

\paragraph{CORPUS REVIEWED.}
The current repository tree was surveyed before new work. Reviewed were \texttt{AGENTS.md}, \texttt{CODEX.md}, the pnp-research skill, root \texttt{main.tex}, TASK-001 and the current task list, newest verified MAIN 1802 and relevant historical MAIN snapshots, LSPV-1 SCHUR1/LEAF2/SU05CODE1/PAIR2COVER1/ORBITSLICE1/SYNCWEDGE1/DIMSUM2/ROOT2ADIC1/HASSETEST1, LSPV-2 WEDGE1/WINDOWREC1/GCD1/DIAGCONV1/SUPPORTK1, newest relevant LSPV-3/LSPV-4 work, governance, product-span code/raw results, and HEC/CRS-PPC/CCSPC/RCCN/Search-G. Repository search covered prior reconstruction-compatible, LDE/RMV, and orbit-label claims. MAIN 1802 remains the newest verified LSPV-5 snapshot and passes substantive/header/sentinel preflight.

\paragraph{Primary-source transcription.}
The 2025 vMS paper fixes the active PCPP branch more explicitly than previously recorded. Lemma 4.2 takes total-degree-$\Delta$ polynomials on $\mathbb F_p^r$, with local-extractor entropy parameter $\Delta^5$, seed length $O(r\log p)$, and output length $\Delta$. Definition 4.3 defines MV lines as axis-parallel combinatorial lines after grouping variables in two halves. RMV enumerates every such line and every local-extractor seed, recursively restricting to those lines.

For the PCPP instantiation, vMS sets
\[
h=m^{100},\quad r=\text{smallest power of two with }h^r\ge |y_i|,
\quad \Delta=hr,
\]
and $p$ to the smallest prime in $[\Delta^{100},2\Delta^{100}]$, and writes $\widehat y_i$ for the canonical low-degree extension. The generator verifies the candidate PCPP $y_i$ with probability one by enumerating all verifier randomness, then outputs $RMV(\widehat y_i)$, along with the encoded-input branch. Thus the relevant polynomial is the canonical LDE of a perfectly accepted PCPP, not an arbitrary low-degree polynomial.

\paragraph{Typing boundary retained.}
vMS Lemma 4.2 imports the local extractor implicitly from SU05 but does not spell out an author-unique internal SU05 matrix, binary code, or field-to-bit label convention. SU05CODE1 remains controlling: SU05 accepts a family of list-decodable binary codes and explicitly permits Hadamard in the computational regime. The bit-level tests below are therefore exact only for the declared source-sanctioned SU05/Hadamard realization, not a claim of a unique vMS implementation.

\paragraph{PROVED/EXACT: canonical LDE leaf typing.}
Place a proof word $y$ on the canonical grid $[h]^r$, padding unused grid locations by zero, and let $\widehat y$ be its canonical tensor-grid LDE. Fixing all but one coordinate to grid values yields a univariate polynomial
\[
f(t)=\widehat y(a_1,\ldots,a_{c-1},t,a_{c+1},\ldots,a_r)
\]
of degree at most $h-1$, whose values on $[h]$ are exactly the corresponding Boolean grid line. LEAF2 proves every such axis-parallel leaf occurs among recursive RMV leaves.

If the padded proof table contains both zero and one, a coordinate path between those grid points has an adjacent change, so at least one axis-parallel Boolean grid line is nonconstant. Thus every nonzero padded proof with at least one zero has a nonconstant legal RMV leaf. Because vMS states only $h^r\ge |y_i|$, strict padding is not automatic in the equality case; the all-one full-grid proof remains a formal degenerate control unless excluded by PCPP structure.

\paragraph{COMPUTATIONAL source-typed leaf battery.}
Runner \texttt{experiments/product_span/vms_lde_hadamard_leaf.py}, commit \texttt{94be4ee3d2de58dd3e6f2a6f057214149ae23c0b}, constructs canonical univariate Boolean-grid LDE leaves over $\mathbb F_p$, evaluates the declared SU05 multiplicative orbit, applies standard binary integer labels plus Hadamard masks, and computes linear and synchronous exterior ranks. Raw results are in \texttt{results/raw/product_span/2026-10-01_LSPV1_vms_lde_hadamard_leaf.json}, result commit \texttt{aadc1efcb9154b67530ac53778fbf4b5fc1a72e5}, seed 314159.

Exhaustive Boolean-grid words were tested at
\[
(p,h,m)=(17,4,4),(29,5,5),(37,6,6),
\]
and 64 deterministic/adversarial/seeded words at $(43,7,7)$. Every tested nonconstant word had
\[
\boxed{\dim W=m,\qquad \dim U=\binom m2.}
\]
No nonconstant exterior-deficient case occurred. Constant controls were deficient. This is \textbf{COMPUTATIONAL}; the small controls do not satisfy the vMS asymptotic $p\ge\Delta^{100}$ and are leaves rather than complete higher-dimensional PCPP instances.

\paragraph{Relation to prior work.}
LEAF2 already observed maximal quadratic feature rank for a larger finite family of interpolated Boolean leaves. The new contribution is the explicit 2025 vMS parameter transcription and a source-typed reimplementation of the leaf test. HASSETEST1 shows repeated-root cyclicity alone does not force low exterior rank. LDELEAF1 shows the same phenomenon persists after imposing canonical Boolean-grid LDE structure in the tested leaf controls.

\paragraph{Literature ideas and status.}
Repeated-root GDFT/Hasse methods remain an exact coordinate transform, not a contraction theorem. Reed--Muller/LDE product theorems motivate source-specific degree filtrations before Boolean labeling, but odd-prime field degree does not bound Boolean ANF degree after labeling. Subspace-design and square-code literature suggests many low-dimensional legal slices can collectively saturate degree two; this is a hostile model, not an SU/vMS theorem. A promising exact route is to use the tensor-grid LDE plus PCPP constraints to classify which univariate grid words can occur on perfectly accepted proofs.

\paragraph{PQ/cross-domain ledger.}
\textbf{EXACT OBJECT/THEOREM TRANSFER:} canonical tensor-grid LDE restriction, RMV axis-parallel leaf enumeration, declared SU05/Hadamard Schur/exterior object, repeated-root Hasse coordinates, Boolean polarization, Gaussian elimination.
\textbf{ANALOGY/CANDIDATE:} derivative-layer shortening, subspace designs, Hilbert/syzygy diagnostics, PCPP-induced leaf constraints, minimum-rank quadratic relations.
\textbf{FAILED TRANSFER:} generic low field degree to low Boolean degree/rank; repeated roots to $O(m)$ exterior rank; generic locality/list-decodability/canonicity to contraction; cryptographic/physical/proof-system analogies as unrestricted SAT runtime.

\paragraph{DUPLICATION / CONTRADICTION AUDIT.}
SCHUR1 already typed the reachable family as LDEs of perfectly accepted PCP proofs. LEAF2 already proved axis-parallel leaf coverage and nonconstant-leaf existence under a padded proof containing both values. ORBITSLICE1 already extracted cyclic windows for a declared Hadamard realization. The new source contribution is the explicit 2025 vMS parameter transcription and source-typed reimplementation. No contradiction was found. The old statement that the grid strictly exceeds proof length is too strong from the rechecked 2025 source: vMS states $h^r\ge |y_i|$, not strict inequality.

\paragraph{WHAT CHANGED RELATIVE TO ROOT MAIN.}
Root MAIN correctly says the actual legal word remains open. We can sharpen this: the PCPP branch fixes canonical LDE parameters and all MV-line recursion, while the bit-level SU05 implementation remains under-specified. For the declared Hadamard realization, finite canonical-LDE leaf controls oppose a generic $O(m)$ exterior contraction. The exact higher-dimensional accepted-PCPP word remains uncomputed.

\paragraph{ADOPTED / REJECTED / RETIRED PRIOR COMPONENTS.}
Adopted: SCHUR1 reachable-family typing, SU05CODE1 encoder boundary, LEAF2 axis-parallel leaf theorem with corrected premise, ORBITSLICE1, SYNCWEDGE1, DIMSUM2, ROOT2ADIC1/HASSETEST1.
Rejected: treating a schematic low-degree polynomial as an exact vMS PCPP instance; treating Hadamard as uniquely canonical.
Retired: generic repeated-root, field-degree, locality, list-decoding, canonicity, or short-source contraction shortcuts.

\paragraph{STALE OR OVERSTATED CLAIM FOUND.}
LEAF2 says vMS chooses $r$ so the grid strictly exceeds proof length. The rechecked 2025 source says $r$ is the smallest power of two such that $h^r\ge |y_i|$. Strict inequality is not stated. The leaf argument remains exact whenever the padded proof contains both values, but the equality/all-one case must not be silently discarded.

\paragraph{Hostile controls and barriers.}
SAME_REALIZATION fixes the accepted PCPP, canonical LDE, legal line, SU05 matrix/code choice, labels, output length and consumer. Semantic/algebraic oracles can reveal all leaves and ranks, so this remains algebraic/relativizing. DEAD/COPY/DELAY preserve extensional tables; INTERP/RECOMP require retyping. Padding cannot be assumed strict. Generic Avoid, affine/Tseitin, Horn, 2-SAT, bounded treewidth, independent products and alternative representations remain controls. Full leaf exterior rank is structural, not runtime hardness.

\paragraph{Strict status and next target.}
\textbf{KNOWN THEOREM/source fact:} vMS canonical-LDE/RMV parameterization and all-line/all-seed enumeration.
\textbf{PROVED/EXACT:} canonical grid restrictions yield Boolean univariate leaves; nonconstant-leaf implication when the padded table contains both zero and one; declared-Hadamard typing.
\textbf{COMPUTATIONAL:} all tested nonconstant source-typed leaves have full $W$ and $U$.
\textbf{FALSIFIED/RETIRED:} generic low-degree/repeated-root structure forces low exterior rank.
\textbf{OPEN:} exclusion of the all-one/full-grid accepted PCPP control; exact higher-dimensional accepted PCPP; unique bit-level vMS realization; asymptotic leaf theorem; global $K_2(H)$; Boolean source degree; charged acquisition; arbitrary-algorithm transfer; nonrelativizing/nonalgebrizing hinge; final diagonal.

The single next target is to inspect the exact PCPP construction used by vMS and decide whether every perfectly accepted proof $y_i$ in the relevant language must contain both a zero and a one, or otherwise has a source-forced nonconstant canonical grid leaf. If yes, combine that source theorem with a symbolic all-parameter proof of the nonconstant-leaf exterior-rank conjecture. If no, construct the accepted constant-proof control explicitly and test whether another legal line or the encoded-input RMV branch supplies the needed exterior span.

\begin{thebibliography}{5}
\bibitem{vMS25} D. van Melkebeek and N. Mocelin Sdroievski, \emph{Instance-Wise Hardness and Refutation versus Derandomization for Arthur--Merlin Protocols}, Computational Complexity 34:15, 2025.
\bibitem{SU05} R. Shaltiel and C. Umans, \emph{Simple Extractors for All Min-Entropies and a New Pseudorandom Generator}, J. ACM 52(2):172--216, 2005.
\bibitem{SU09} R. Shaltiel and C. Umans, \emph{Low-End Uniform Hardness versus Randomness Trade-offs for AM}, SIAM J. Comput. 39(3):1006--1037, 2009.
\bibitem{Ling} S. Ling, H. Niederreiter, P. Sol\'e, \emph{On the Algebraic Structure of Quasi-cyclic Codes IV: Repeated Roots}, Designs Codes Cryptogr. 38:337--361, 2006.
\bibitem{BGHSV} E. Ben-Sasson, O. Goldreich, P. Harsha, M. Sudan, S. Vadhan, \emph{Short PCPs Verifiable in Polylogarithmic Time}, CCC 2005.
\end{thebibliography}
\end{document}
% END_OF_VERIFIED_REPORT::LSPV-1::LDELEAF1
